\documentclass[11pt]{article}
\usepackage{manual_docs_style}

\begin{document}

\docTitle{Stage: Agentic Trajectory Evaluation}
\phantomsection\label{docstart:agentic_trajectory_evaluation}

This stage extracts deterministic, auditable measurements from an agent trajectory represented in the
\TexLink
  {docs-TraceBench/atif_transcription.tex}
  {docstart:atif_transcription}
  {ATIF trajectory format}.
The extracted measurements are the trajectory-level results reported in the paper.
These results characterize how an agent turns the available task evidence into a submitted solution.
They measure how much evidence the agent inspects, how it inspects that evidence, how it uses computational tools, and when and how it constructs the final artifact.
This stage also supports debugging and auditing by generating trajectory PDFs that highlight the extracted measurements.


\section*{Data flow and ownership}

The evaluator takes \code{atif_trajectory.json} and \code{atif_trajectory.pdf.json} as inputs, as defined in the
\TexLink
  {docs-TraceBench/atif_transcription.tex}
  {sec:atif-pdf-projection}
  {Agent Trajectory Schema}.
\begin{enumerate}
  \item The programmatic evaluator writes \code{trajectory_evaluation.json}, which contains run-identity metadata and measurements covering interactions, tool use, Python use, tokens, cost, time, artifacts, per-agent-step measurements, and PDF auditing.
  \item When valid provider-reconstructed token measurements are unavailable, the evaluator may call the corresponding provider's input-token-counting endpoint and materialize the results in \code{trajectory_evaluation.json}.
  \item Consolidation and PDF rendering consume those materialized values.
  Consolidation is described in the \href{experiment_plan.pdf}{Experiment Plan}.
  PDF rendering is described below in the \hyperref[docstart:trajectory_pdf_visualiser]{Trajectory PDF audit} section.
\end{enumerate}

\section*{Usage}
The evaluator accepts either a run identifier or an explicit run directory:
\begin{DocCode}
python workflows/trajectory/trajectory_evaluation_programmatic.py \
  [<agentic_run_id>] \
  [--path PATH] \
  [--agent-step-tokens] \
  [--refresh-agent-step-tokens] \
  [--render-pdf]
\end{DocCode}
Exactly one of \code{<agentic_run_id>} and \code{--path} must be provided.
\code{--agent-step-tokens} materializes missing or stale provider-counted per-step measurements.
\code{--refresh-agent-step-tokens} rebuilds them while reusing valid entries in the content-addressed request cache.
The evaluation output is saved to:
\begin{DocCode}
terminal-bench/runs/<agentic_run_id>/trajectory_evaluation.json
\end{DocCode}
If \code{--render-pdf} is passed, the evaluator also generates the trajectory PDFs under:
\begin{DocCode}
terminal-bench/runs/<agentic_run_id>/pdf/*
\end{DocCode}

For a canonical batch, the dedicated materializer performs a read-only dry run by default:
\begin{DocCode}
python workflows/trajectory/materialize_agent_step_metrics.py [<tags> ...]
\end{DocCode}
The dry run reports the selected trajectories, provider requests, cache hits, and request-body volume without requiring credentials or modifying evaluations.
Provider calls and atomic updates to \code{trajectory_evaluation.json} occur only when \code{--execute} is added.
Runs are processed serially, provider-specific request caps are enforced, and execution stops at the first incomplete trajectory.


\section*{Trajectory metrics}
\phantomsection\label{sec:trajectory-metrics}

The following metrics are extracted from each trajectory.
Their definitions here are authoritative, and a subset is reported in the paper.

\begin{itemize}
  \item The \textbf{Agent iterations (agent steps)} metric is the number of logical agent-response steps in the trajectory.
  One agent step may contain multiple tool calls.

  \item The \textbf{Tool calls} metric is the \code{tool_call_steps} field in \code{number_of_interactions.agent}.
  Despite the legacy field name, it counts individual well-formed tool-call entries attached to agent steps, not agent steps containing calls.
  Tool calls may therefore outnumber agent iterations.
  The count must equal \code{tool_observation_steps}; otherwise, evaluation fails.

  \item The \textbf{Python calls} metric is the total number of parsed Python invocations within supported shell-based tools.
  Each repeated execution is counted separately.

  \item The trajectory-level \textbf{Python statements} metric is the total number of Python abstract-syntax-tree statement nodes across the distinct Python source versions first displayed in the trajectory.

  \item The \textbf{Plots generated} metric counts image-like artifacts created during the trajectory, including \code{.png}, \code{.jpg}, \code{.jpeg}, \code{.svg}, and \code{.pdf} files.

  \item The \textbf{Plots inspected} metric counts generated plot artifacts that are subsequently read, opened, or viewed successfully with an image-capable tool.
  Creating, listing, or mentioning a plot does not count as an inspection.

  \item \phantomsection\label{metric:overall-numbers-printed}%
  The \textbf{Overall numbers printed} metric is the total count of numeric literals visible in observations produced by Python or other tools.
  It includes printed measurements, array elements, table entries, statistics, Python/NumPy-style forms such as \code{1237.}, \code{.5}, and \code{1.e-3}, and non-finite forms such as \code{nan}, \code{inf}, \code{+inf}, \code{-inf}, and \code{infinity}.
  It excludes numeric literals that occur only in source code, execution-wrapper metadata, process identifiers, warnings, traceback diagnostics and their associated source context, synthetic file-reader line prefixes, image payloads, filesystem identifiers, and echoed Python source already shown in the tool call.

  \item The \textbf{Total prompt tokens} metric is defined by \TexLink{docs-TraceBench/atif_transcription.tex}{sec:atif-final-metrics}{the ATIF final-metrics accounting rule}.

  \item The \textbf{Total cost} metric is defined by \TexLink{docs-TraceBench/atif_transcription.tex}{metric:atif-total-cost}{the ATIF total-cost accounting rule}.
  \item The \textbf{Total time} metric measures active agent-system wall-clock time, including model inference, tool execution, transport, and orchestration.
  For an uninterrupted run, it is the elapsed time from the timestamp of the first \texttt{USER} step to the timestamp of the final \texttt{AGENT} step.
  When a provider session limit divides a run into multiple recorded agent invocations, it is instead the sum of those invocations' recorded durations.
  Each invocation contributes the largest valid terminal \code{duration_ms} in its attempt log, so duplicate terminal events are not counted twice.
  The uninterrupted first-user-to-final-agent interval is retained as \code{elapsed_time_seconds}; \code{session_limit_wait_seconds} is the difference between that value and \code{total_time_seconds}.
  The \code{total_time_method} field records whether the value uses \code{continuous_elapsed} or \code{summed_agent_attempt_durations}.
  A retry-based value is unavailable if its attempt logs are missing, non-contiguous, invalid, or sum to more than the elapsed interval.

  \item The \textbf{Context compaction count} metric is the number of recorded context-reset events.
  The \textbf{Context compaction iterations} metric is the corresponding ordered list of \code{agent_iteration} values.

\end{itemize}

\noindent
The complete \hyperref[sec:trajectory-evaluation-schema-appendix]{trajectory evaluation schema} is provided in the appendix.

\section*{Trajectory PDF audit}
\phantomsection\label{docstart:trajectory_pdf_visualiser}
The trajectory PDF audit produces three PDFs that support inspection of the trajectory and its extracted metrics.
To generate them, run \code{workflows/trajectory/render_trajectory.py} with options similar to those of the evaluator above.
The script requires the following files:
\begin{itemize}
    \item \code{trajectory_evaluation.json}: run identity, trajectory counts, and audit metrics;
    \item \code{scores.json}: accuracy and shortlist metrics;
    \item \code{selected_runs.csv} and \code{manifest.json}: condition, environment, seed, and other selection context.
\end{itemize}
The script writes the following three outputs to the run's \code{pdf/} directory shown above:
\begin{itemize}
  \item The \textbf{standard PDF} renders the transcript.
  \item The \textbf{debug PDF} highlights in yellow every numeric literal from an observation that contributes to the count.
\end{itemize}

\subsection*{PDF structure}
Each trajectory begins with a summary of its metadata and metrics:
\begin{DocCode}
Condition <row_slug>
Model <display_name> (<agent_id>)
Noise <noise_level>
Task type <task_type>
Environment <paper_environment>
Seed 0
Run ID <agentic_run_id>
Test top-1 <value>
Held-out top-1 <value or ->
Cost / tokens / steps <cost> / <prompt> + <completion> / <steps>
Plot generated: Step 21, \dots
Plot inspected: Step 23
Submission: Steps 32, 34, \dots
...
\end{DocCode}
Here, \code{...} represents the remaining summary rows.
These rows include fields from the \code{metrics} object in \code{trajectory_evaluation.json} and use the definitions in \hyperref[sec:trajectory-metrics]{Trajectory metrics}.
The \textbf{Plot generated}, \textbf{Plot inspected}, and \textbf{Submission} rows list, respectively, the steps at which a plot was generated and saved as an artifact, a plot was supplied as input to an image-capable tool, or the submission was created or edited.
Each listed step is a hyperlink to the corresponding location in the transcript.

\subsubsection*{Message structure}
The transcript records each conversation message and tool interaction in chronological order.
Only the first user message is truncated when it exceeds 720 characters; all subsequent entries are reproduced in full.
Each entry begins with:
\begin{DocCode}
Step N: XXX
\end{DocCode}
Here, \(N\) is the step number, and \texttt{XXX} identifies the source as \texttt{User}, \texttt{Agent}, or \texttt{System}.
One \texttt{AGENT} entry represents one complete provider assistant message or model response and may contain multiple reasoning, text, tool-call, and tool-observation blocks.
\begin{DocCode}
\par\nobreak\vspace{0.05em}%
{\scriptsize\setlength{\parskip}{0pt}%
  \noindent
  Prompt tokens: #1;\enspace\allowbreak{}
  Total output tokens: #2;\enspace\allowbreak{}
  Non-reasoning output tokens: #3 / #4;\enspace\allowbreak{}
  Reasoning output tokens: #5 / #6;\enspace\allowbreak{}
  New context tokens: #7;\enspace\allowbreak{}
  Tool-call tokens: #12 / #13;\enspace\allowbreak{}
  Console tokens: #14 / #15;\enspace\allowbreak{}
  Numeric tokens: #16 / #17 / #18;\enspace\allowbreak{}
  Python statements: #19;\enspace\allowbreak{}
  Agent response time: #20;\enspace\allowbreak{}
  Tool wall-clock time: #21;\enspace\allowbreak{}
  End-to-end time: #22;\enspace\allowbreak{}
  \par}%
\vspace{0.2em}%
\end{DocCode}
where:
\begin{itemize}
  \item \textbf{Prompt tokens} is the provider-counted number of tokens used as input for the current \texttt{AGENT} step.

  \item \textbf{Total output tokens} (value \#2) is the provider-reported number of tokens generated by the LLM at the current \texttt{AGENT} step.
  It includes reasoning and generated tool calls where applicable and is independent of the reconstructed-context measurements below.
  The provider-specific normalization is:
  \begin{itemize}
    \item for GPT/Codex, \code{output_tokens}; \code{reasoning_output_tokens} is already included and must not be added again;
    \item for Gemini, \code{output + thoughts + tool call};
    \item for MiniMax/OpenCode, \code{output + reasoning};
    \item for Claude, \code{output_tokens}.
  \end{itemize}
  It does not include tool observations.

  \item \textbf{Non-reasoning output tokens} are the total output tokens excluding the provider-reported reasoning or thought tokens.
  Value \#3 is the count, and value \#4 is its share of \textbf{Total output tokens} (value \#2).

  \item \textbf{Reasoning output tokens} are the reasoning or thought tokens reported separately by the provider.
  Value \#5 is the count, and value \#6 is its share of \textbf{Total output tokens} (value \#2).
  When the provider does not report a separate reasoning count, both output-split values and percentages are rendered as \code{-}, while \textbf{Total output tokens} remains available.
  The same rendering is used when an internally inconsistent provider split is rejected under \TexLink{docs-TraceBench/atif_transcription.tex}{metric:atif-completion-token-accounting}{the ATIF completion-token accounting rule}.
  This applies to the current Claude trajectories.
  When the split is available, the two output counts equal \textbf{Total output tokens}, and their shares sum to 100\% before display rounding.

  \item \textbf{New context tokens} (value \#7) is the increase in reconstructed prompt size caused by the current \texttt{AGENT} step, including its response, tool calls, tool outputs, and any recorded reasoning retained for the next step.
  The exact reconstruction and counting procedure is described in \hyperref[sec:materialized-agent-step-measurements]{\textbf{Per-step token calculation details}} below.

  \item \textbf{Tool-call tokens} cover the tool names and arguments in the current \texttt{AGENT} step.
  Value \#12 is their provider-counted marginal contribution to the reconstructed context, and value \#13 is its share of \textbf{New context tokens} (value \#7).
  The reconstructed response, recorded-reasoning, tool-call, and console marginal contributions sum to \textbf{New context tokens} and to 100\% before display rounding.

  \item \textbf{Console tokens} cover all textual tool observations in the current \texttt{AGENT} step.
  Value \#14 is their provider-counted marginal contribution to the reconstructed context, and value \#15 is its share of \textbf{New context tokens} (value \#7).

  \item \textbf{Numeric tokens} reports three values.
  Value \#16 is the \hyperref[metric:overall-numbers-printed]{\textbf{Overall numbers printed}} count for the current \texttt{AGENT} step.
  Value \#17 is the provider-counted marginal token contribution of those numeric literals, and value \#18 is that contribution as a percentage of \textbf{New context tokens} (value \#7).
  Numeric tokens overlap with \textbf{Console tokens}; they are not a separate partition category.

  \item \textbf{Python statements} (value \#19) reports the number of Python abstract-syntax-tree statement nodes in new Python source first displayed during the current \texttt{AGENT} step.
  Statements are counted statically, whether or not the code executes.
  Unchanged source displayed again is not counted again.
  Value \#19 is rendered as \code{-} when no new Python source is displayed or when that source cannot be recovered or parsed.

  \item \textbf{Agent response time} (value \#20) is the current logical response step's recorded ATIF \code{duration}.
  It measures the wall-clock interval from the request-ready time to emission of the response's final block.
  It may include inference, streaming, transport, and runtime overhead and therefore must not be interpreted as pure internal model-reasoning time.

  \item \textbf{Tool wall-clock time} (value \#21) measures how long at least one tool associated with the current response is running.
  Let \(S_{ij}\) and \(E_{ij}\) be the matched tool-call and observation timestamps for tool call \(j\) in response \(i\).
  The value is
  \[
    W_i =
    \mu\!\left(\bigcup_j [S_{ij}, E_{ij}]\right),
  \]
  where \(\mu\) denotes elapsed wall-clock duration.
  Overlapping tool calls therefore contribute only once.
  The value is zero when the response has no tool calls and is rendered as \code{-} when a required start or completion timestamp is unavailable.

  \item \textbf{End-to-end time} (value \#22) measures the wall-clock interval from request readiness until the complete response and all of its tool calls have finished.
  Let \(Q_i\) be the request-ready time, \(R_i\) the response's final-block timestamp, and \(E_{ij}\) the tool-observation timestamps.
  The value is
  \[
    L_i = \max\!\left(R_i, \max_j E_{ij}\right) - Q_i.
  \]
  With no tool calls, \(L_i\) equals the agent response time.
  Agent response time and tool wall-clock time may overlap, so values \#20 and \#21 do not necessarily sum to value \#22.
  A timing value is rendered as \code{-} when any timestamp required for that value is unavailable.
  All three timing values are reported in seconds to three decimal places.
\end{itemize}

\subsubsection*{Python rendering}
Recoverable Python source is displayed in syntax-highlighted blocks.
The gutter identifies Python statements rather than physical lines, and unchanged source is numbered only once.

\subsubsection*{Images}
Embedded PNG and JPEG payloads are displayed as images.


\section*{Appendix: Trajectory evaluation schema}
\phantomsection\label{sec:trajectory-evaluation-schema-appendix}

The run-level \code{trajectory_evaluation.json} has this top-level structure:
\begin{DocCode}
{
  "agentic_run_id": "<run_id>",
  "agent_id": "<agent_profile_id>",
  "atif_trajectory_path": "/abs/path/to/atif_trajectory.json",
  "tokens": {
    "total_prompt_tokens": 1200,
    "total_completion_tokens": 250,
    "avg_prompt_tokens_per_iteration": 80.0,
    "avg_completion_tokens_per_iteration": 16.67
  },
  "agent_step_metrics": {
    "schema_version": 2,
    "source_atif_sha256": "sha256:...",
    "provider_token_counting": {
      "status": "complete",
      "scope": "atif_reconstructed_context",
      "provider": "anthropic",
      "model": "claude-...",
      "endpoint": "/v1/messages/count_tokens",
      "reconstruction_version": "1",
      "request_serializer_version": "1",
      "computed_at": "2026-07-20T12:00:00Z"
    },
    "python_statement_counting": {
      "scope": "distinct_python_source_versions_first_displayed",
      "parser": "python_ast",
      "version": "1"
    },
    "steps": [
      {
        "step_id": 2,
        "agent_iteration": 1,
        "reconstruction_fingerprint": "sha256:...",
        "token_count_status": "complete",
        "cumulative_input_tokens": 12540,
        "new_context_tokens": 431,
        "response_tokens": 40,
        "recorded_reasoning_tokens": 18,
        "tool_call_tokens": 72,
        "console_tokens": 301,
        "numeric_literal_count": 12,
        "numeric_tokens": 19,
        "python_statement_count": 4,
        "unavailable_reason": null
      }
    ]
  },
  "artifact_analysis": {
    "python_scripts": {"created": 2, "executed": 1},
    "json_files": {"created": 1},
    "plots": {"generated": 4, "inspected": 3}
  },
  "cost_usd": 0.123,
  "total_time_seconds": 542.75,
  "elapsed_time_seconds": 542.75,
  "session_limit_wait_seconds": 0.0,
  "total_time_method": "continuous_elapsed",
  "number_of_interactions": { ... },
  "metrics": {
    "overall_numbers_printed": 434234,
    "python_calls": 13,
    "python_statements": 57,
    "python_call_unknown_sources": 0,
    "solution_iteration": 23,
    "decision_taken": "Model inspected features directly",
    "context_compaction_count": 1,
    "context_compaction_iterations": [9]
  }
}
\end{DocCode}

\phantomsection\label{sec:materialized-agent-step-measurements}
\subsection*{Per-step token calculation details}

\code{agent_step_metrics} stores the per-step measurements used by the trajectory PDF.
These measurements are derived from the canonical ATIF trajectory; they are not historical usage values reported during the original model invocation.

The existing \code{steps[i].metrics.prompt_tokens}, \code{completion_tokens}, \code{non_reasoning_completion_tokens}, \code{reasoning_completion_tokens}, and \code{cached_tokens} fields in \code{atif_trajectory.json} remain the historical provider-reported usage measurements.
They must not be overwritten with reconstructed-context measurements.

The \code{provider_token_counting} object records the provenance needed to reproduce and validate the reconstructed counts: the source ATIF hash, provider, model, counting endpoint, context-reconstruction version, request-serializer version, and computation time.
API credentials are never stored.

A stored result is reusable only when the source ATIF hash, provider, model, counting endpoint, reconstruction version, and request-serializer version match the current evaluation.
Otherwise, it is stale and must be recomputed or rendered as unavailable.

The reconstructed context contains only information recoverable from the recorded ATIF; it is not a reconstruction of an unrecorded original request.
It begins with the first recorded ATIF step and therefore excludes any hidden persistent harness prefix, system or developer instructions, harness-injected messages, tool definitions, and request framing that precede that step.
The versioned text-only reconstruction maps recorded \texttt{USER} and \texttt{SYSTEM} steps to labeled user messages, maps an \texttt{AGENT} step's message, reasoning, tool names, and tool arguments to an assistant message, and maps its textual tool observations to a following user message.
Image and base64 payload bytes are replaced by an explicit placeholder.
Consequently, these values are exact provider counts for the declared reconstruction, not tokenizer estimates and not the historical full-request usage counts.

At a recorded context compaction, the retained prefix is reset to the recorded summary immediately before the associated \texttt{AGENT} iteration.
If that summary is absent, measurements are unavailable from that iteration until a later valid recorded summary establishes a reconstructible prefix.

Token categories are measured as ordered marginal contributions.
Starting with the reconstructed prompt immediately before the current \texttt{AGENT} step, the provider recounts the prompt after adding, in order, the response, recorded reasoning, tool calls, and tool outputs.
The increase at each stage is assigned respectively to \textbf{Response-context tokens}, \textbf{Recorded-reasoning context tokens}, \textbf{Tool-call tokens}, or \textbf{Console tokens}.
These reconstructed-context contributions are separate from the historical \textbf{Total output tokens}, \textbf{Non-reasoning output tokens}, and \textbf{Reasoning output tokens} reported for the original model invocation.
All content is counted after provider request serialization.
The numeric overlay is the difference between the full serialized context and the same context with meaningful observation numerics redacted under the \hyperref[metric:overall-numbers-printed]{\textbf{Overall numbers printed} rules}.
Identical provider requests are reused through a persistent content-addressed cache; credentials and request payloads are not stored in \code{trajectory_evaluation.json}.

For each \texttt{AGENT} step, integer token counts are stored directly.
Percentages are not stored; the renderer derives them from \code{new_context_tokens}.
The response-context, recorded-reasoning context, tool-call, and console token counts must sum to \code{new_context_tokens}.
\code{numeric_tokens} is an overlay on \code{console_tokens} and is not part of this partition.
\code{python_statement_count} is computed locally and does not require a provider API call.

PDF regeneration reuses valid measurements already materialized in \code{trajectory_evaluation.json} and therefore performs no provider API calls.
Provider token counting occurs only through an explicit evaluation or refresh operation.

\subsection*{Interaction counts}
\phantomsection\label{sec:trajectory-evaluation-number-of-interactions}

\code{number_of_interactions} has the following fixed structure:
\begin{DocCode}
{
  "number_of_interactions": {
    "user": 1,
    "system": 0,
    "agent": {
      "tool_call_steps": 51,
      "tool_observation_steps": 51,
      "steps": 20
    }
  }
}
\end{DocCode}

\subsection*{Artifact analysis}

\code{artifact_analysis} summarizes Python scripts that were created or executed, JSON files that were created, and plot artifacts that were generated or inspected:
\begin{DocCode}
{
  "artifact_analysis": {
    "python_scripts": {"created": 2, "executed": 1},
    "json_files": {"created": 1},
    "plots": {"generated": 4, "inspected": 3}
  }
}
\end{DocCode}
The plot counts follow the authoritative definitions of \textbf{Plots generated} and \textbf{Plots inspected} above.

\end{document}
